\chapter{Propuesta de Solución}

A partir de los hallazgos obtenidos en la revisión sistemática de la literatura, esta sección presenta la propuesta de solución definida para abordar las principales brechas identificadas en el análisis de logs de seguridad mediante modelos de lenguaje de gran escala. La propuesta consiste en el diseño de una arquitectura prototipo, orientada al análisis semántico y agentivo de logs de seguridad mediante modelos de lenguaje de gran escala. La solución se concibe como un sistema de apoyo al analista de incidentes, capaz de procesar eventos provenientes de diversas fuentes, extraer información relevante, formular hipótesis sobre posibles amenazas y generar una reconstrucción estructurada del incidente observado. Actualmente, se cuenta con una versión preliminar del agente, por lo que esta sección no solo describe la arquitectura conceptual propuesta, sino también el trabajo desarrollado hasta esta etapa del seminario.

\section{Descripción General}

A diferencia de los enfoques tradicionales basados exclusivamente en reglas, clasificación automática o generación de resúmenes, la arquitectura propuesta incorpora un flujo de razonamiento por fases. En dicho flujo, distintos módulos especializados cumplen funciones complementarias, entre ellas el preprocesamiento de logs, el análisis inicial del evento, la recuperación de conocimiento, la consulta de fuentes externas, la generación de mitigaciones, la validación de acciones y la elaboración de un reporte final. Así, el sistema no se limita a entregar una predicción aislada, sino que organiza el análisis como una secuencia controlada de pasos orientados a fundamentar el diagnóstico.


La estructura general de la propuesta se presenta en la Figura~\ref{fig:diag_general}. El primer componente de la solución corresponde al preprocesamiento determinístico de los logs. Esta etapa tiene como propósito transformar registros crudos en una representación estructurada, identificando elementos como direcciones IP, usuarios, servicios afectados, eventos críticos, frecuencia de ocurrencia y ventanas temporales. Esta transformación permite reducir el ruido de entrada, facilitar el análisis posterior y evitar que el modelo de lenguaje opere directamente sobre grandes volúmenes de información no depurada.

Luego, la arquitectura incorpora un módulo de razonamiento basado en agentes, encargado de interpretar los eventos procesados y determinar si la evidencia disponible es suficiente para emitir un diagnóstico preliminar. Cuando se requiere mayor contexto, el sistema puede activar mecanismos de recuperación aumentada por conocimiento, consultar una memoria vectorial o acceder a fuentes externas de inteligencia de amenazas. Esta base de conocimiento puede ser poblada con información estructurada de MITRE ATT\&CK, lo que permitiría asociar los eventos detectados con tácticas, técnicas, subtécnicas, procedimientos y posibles grupos de amenaza.

La solución también contempla un módulo de generación y validación de respuestas. Este componente produce acciones de mitigación, incorporando una revisión previa destinada a identificar recomendaciones inseguras, ambiguas o potencialmente disruptivas. El proceso culmina con la generación de un informe estructurado que sintetiza la hipótesis del incidente, la evidencia utilizada, las relaciones causales inferidas, las técnicas asociadas y el nivel de confianza del análisis.


\begin{figure}[H]
    \centering
    \includegraphics[width=0.8\textwidth]{Resource/Diagrama_general.png}
    \caption{Daigrama general de procesamiento para el análisis de logs de seguridad mediante un agente basado en LLM. Fuente: Elaboración propia.}
    \label{fig:diag_general}
\end{figure}

%%-------------------------------------------------------------
\section{Arquitectura general del sistema}
%%-------------------------------------------------------------

La arquitectura propuesta organiza el análisis de logs de
seguridad como un \textit{pipeline Plan-and-Execute}, que
separa explícitamente la fase de planificación de la fase de
ejecución. Esta decisión de diseño responde a una limitación
documentada del patrón ReAct tradicional: su carácter miope,
donde el agente decide el siguiente paso únicamente con lo
observado hasta ese momento, sin una visión global del
problema~\cite{rigaki2023cage}. En Plan-and-Execute, un módulo
Planificador construye primero un plan completo sobre un resumen
del problema, y un módulo Ejecutor lo implementa paso a paso
mediante \textit{function calling} nativo con el modelo de
lenguaje.

El pipeline completo se estructura en cuatro etapas
secuenciales, ilustradas en la Figura~\ref{fig:arq_general}:
(1)~agregación determinística de los logs de entrada,
(2)~planificación del análisis mediante un agente LLM,
(3)~ejecución agéntica a través de un grafo de estados
construido con LangGraph, y (4)~síntesis y generación del
reporte final. La Figura~\ref{fig:arq_general} resume el flujo
completo del sistema, desde la ingesta de logs crudos hasta la
entrega del reporte estructurado al analista.

\begin{figure}[H]
    \centering
    \includegraphics[width=0.9\textwidth]{Resource/plan_exc.png}
    \caption{Arquitectura general del sistema propuesto.
             El flujo procede desde la ingesta de logs crudos
             hasta la generación del reporte, pasando por las
             etapas de agregación determinística, planificación,
             ejecución agéntica con LangGraph y síntesis.
             Fuente: Elaboración propia.}
    \label{fig:arq_general}
\end{figure}

%%-------------------------------------------------------------
\subsection{Preprocesamiento y agregación de logs}
%%-------------------------------------------------------------

La primera etapa recibe el lote de logs crudos en formatos
heterogéneos (\textit{e.g.},sshd, Apache, kernel e IDS) y produce un resumen
estructurado sin invocar al modelo de lenguaje. Esta
decisión es deliberada: enviar directamente miles de líneas de
log a un LLM es costoso, lento y propenso a errores de
interpretación~\cite{liu2024logprompt}. En su lugar, un parser
determinístico basado en expresiones regulares extrae por cada
línea la dirección IP de origen, el usuario, el servicio
afectado y el tipo de evento. Una función de agregación
consolida el lote en un resumen con las siguientes dimensiones:
total de líneas procesadas, número de IPs distintas,
distribución de servicios, distribución de eventos, ventana
temporal y estadísticas por IP (intentos fallidos, número de
eventos).

En paralelo, todas las IPs distintas identificadas en el lote
se corroboran contra la API de AbuseIPDB mediante un módulo de
caché persistente. Cada IP se consulta exactamente una vez:
las consultas posteriores se sirven desde un archivo de caché
local, evitando el uso excesivo de la API y reduciendo la
latencia. Cada IP recibe una categoría según umbrales
configurables: \textit{malicioso} (score~$\geq 75$),
\textit{sospechoso} (score~$\geq 25$) o \textit{limpio}
(score~$< 25$). El resultado de esta etapa es un resumen
estructurado que actúa como único input de los módulos de
razonamiento posteriores, lo que garantiza que el modelo de
lenguaje opere sobre evidencia depurada y compacta en lugar
de sobre registros crudos.

\begin{algorithm}[H]
\caption{Preprocesamiento y agregación de logs}
\label{alg:agregacion}
\KwIn{Lote de logs crudos $L$ (sshd, Apache, kernel, IDS)}
\KwOut{Resumen estructurado $S$, mapa de reputación $R$}
\ForEach{línea $l \in L$}{
    Extraer IP, usuario, servicio y tipo de evento mediante regex\;
    Acumular estadísticas por IP y por servicio\;
}
Construir resumen $S$ con total de líneas, IPs distintas,
distribución de servicios, eventos y ventana temporal\;
\ForEach{IP $ip$ en $S.\text{ips\_distintas}$}{
    \eIf{$ip \in \text{caché}$}{
        $R[ip] \leftarrow \text{caché}[ip]$\;
    }{
        Consultar AbuseIPDB; categorizar por score\;
        $R[ip] \leftarrow \text{resultado}$; actualizar caché\;
    }
}
\Return{$S$, $R$}\;
\end{algorithm}


%%-------------------------------------------------------------
\subsection{Módulo de planificación}
%%-------------------------------------------------------------

El módulo Planificador recibe el resumen estructurado $S$ y
produce un plan ordenado de pasos antes de ejecutar ninguna
acción. Su función es descomponer el problema de análisis en
una secuencia de herramientas a invocar, con un objetivo
explícito para cada paso. Para enriquecer el plan, el
Planificador consulta la memoria de largo plazo (ChromaDB):
recupera hasta $k=3$ patrones de ataque similares para
\textit{few-shot} dinámico, y hasta $k=3$ correcciones
históricas de errores pasados, ambos recuperados por similitud
coseno con un umbral mínimo de 0.55. Este mecanismo permite
que el agente aproveche el conocimiento acumulado en análisis
anteriores, reduciendo la probabilidad de repetir errores
previos~\cite{lee2025logrespagent}.

El plan resultante especifica, para cada paso, la herramienta
a invocar y el objetivo concreto, con un máximo configurable
de 6 pasos. Si la ejecución del plan no produce evidencia
útil, el sistema activa un mecanismo de re-planificación con
un enfoque alternativo, permitiendo hasta 1 re-planificación
por análisis. El Algoritmo~\ref{alg:planificacion} describe
el proceso de planificación y re-planificación.

Cabe señalar que los parámetros de recuperación semántica
($k=3$, umbral de similitud 0.55) están calibrados para el
volumen experimental del prototipo actual, cuyas colecciones
ChromaDB contienen un número reducido de entradas. En el futuro Proyecto de Título se evaluará el ajuste de estos parámetros
al operar con bases de conocimiento de mayor escala y
diversidad temática.

\begin{algorithm}[H]
\caption{Módulo de planificación}
\label{alg:planificacion}
\KwIn{Resumen $S$, mapa de reputación $R$}
\KwOut{Plan de pasos $P$}
Recuperar patrones similares $E \leftarrow
    \text{ChromaDB.retrieve}(S.\text{firma}, k=3)$\;
Recuperar correcciones $C \leftarrow
    \text{ChromaDB.errors}(S.\text{firma}, k=3)$\;
Construir prompt con $S$, $R$, $E$ y $C$\;
$P \leftarrow \text{LLM.plan}(\text{prompt})$
    \tcp*{máx. 6 pasos}
\eIf{$P.\text{pasos} = \emptyset$}{
    $P \leftarrow \text{LLM.replan}(S)$
        \tcp*{plan alternativo, máx. 1 vez}
}{
    \Return{$P$}\;
}
\Return{$P$}\;
\end{algorithm}

%%-------------------------------------------------------------
\subsection{Módulo de ejecución agéntica con LangGraph}
%%-------------------------------------------------------------

La ejecución se implementa como una máquina de estados formal construida con LangGraph, compuesta por 7~nodos (4~agentes LLM y 3~herramientas determinísticas) y 3~aristas condicionales denominadas \textit{routers}. Esta estructura permite representar el análisis como un flujo de decisiones condicionales, donde el grafo determina dinámicamente qué nodo se activa a continuación según el estado actual del sistema~\cite{blefari2026cyberrag}. La representación general de este flujo se puede observar en la Figura~\ref{fig:componentes}.


El flujo de ejecución es el siguiente:

\begin{enumerate}
    \item El \textbf{Agente Triage} analiza el resumen del lote y decide si requiere
     contexto histórico de ChromaDB; si no dispone de contexto,
    activa el router hacia la herramienta RAG.
    
    \item La \textbf{Herramienta RAG} recupera patrones similares de la colección \texttt{attack\_patterns} y los inyecta como contexto al Agente Triage, que genera entonces un diagnóstico preliminar.
    
    \item El \textbf{Agente Selector de Herramientas}
     evalúa el diagnóstico y decide si es necesario consultar
    vulnerabilidades CVE en NVD para los servicios afectados
    
    \item La \textbf{Herramienta CVE} consulta la API de NVD y retorna
    los CVEs relevantes
   
    \item El \textbf{Agente Ingeniero de
    Mitigación} genera comandos de mitigación Linux (iptables,
    fail2ban, ufw, systemctl); si los comandos son rechazados por
    el validador, el agente los corrige en un ciclo de
    autorreflexión con un máximo de 2~intentos.
   
    \item La \textbf{Herramienta Validador} verifica los comandos contra
    una lista blanca de patrones seguros, rechazando con motivo
  explícito aquellos que no coincidan.
    
    \item El \textbf{Agente Auditor Final} evalúa la calidad global del análisis y asigna
     un nivel de confianza $c \in [0,1]$ al resultado.
\end{enumerate}

\begin{figure}[H]
    \centering
    \includegraphics[width=0.43\textwidth]{Resource/componentes.png}
    \caption{Grafo de orquestación implementado mediante LangGraph para la versión preliminar del agente.Fuente: Elaboración propia.}
    \label{fig:componentes}
\end{figure}

El Algoritmo~\ref{alg:ejecucion} resume el flujo del grafo
LangGraph, incluyendo los tres routers condicionales y el
ciclo de corrección de comandos.

\begin{algorithm}[H]
\caption{Ejecución agéntica con LangGraph}
\label{alg:ejecucion}
\KwIn{Resumen $S$, plan $P$, mapa de reputación $R$}
\KwOut{Diagnóstico $D$, comandos validados $C_v$,
        confianza $c$}
\tcp{Nodo 1: Triage}
\eIf{ChromaDB contiene patrones similares}{
    $D \leftarrow \text{Triage.LLM}(S, R,
        \text{contexto\_RAG})$\;
}{
    contexto $\leftarrow \text{RAG.ChromaDB}(S.\text{firma})$\;
    $D \leftarrow \text{Triage.LLM}(S, R, \text{contexto})$\;
}
\tcp{Nodo 3: Selector de herramientas}
\eIf{$D$ menciona servicio con posibles CVEs}{
    $\text{CVEs} \leftarrow \text{NVD.query}(D.\text{servicio})$\;
}{
    $\text{CVEs} \leftarrow \emptyset$\;
}
\tcp{Nodo 5--6: Mitigación y validación (ciclo)}
$\text{intentos} \leftarrow 0$\;
\Repeat{comandos válidos \textbf{o} $\text{intentos} \geq 2$}{
    $C \leftarrow \text{Mitigacion.LLM}(D, \text{CVEs},
        \text{motivos\_rechazo})$\;
    $(v, \text{motivos}) \leftarrow
        \text{Validador}(C)$\;
    $\text{intentos} \leftarrow \text{intentos} + 1$\;
}
$C_v \leftarrow C$\;
\tcp{Nodo 7: Auditor final}
$c \leftarrow \text{Auditor.LLM}(D, C_v, \text{CVEs})$\;
\Return{$D$, $C_v$, $c$}\;
\end{algorithm}

%%-------------------------------------------------------------
\subsection{Síntesis y generación del reporte}
%%-------------------------------------------------------------

El módulo de síntesis integra el resumen del lote $S$, la
reputación de las IPs corroboradas $R$ y las observaciones
de las herramientas para producir el diagnóstico global.

Si la confianza del análisis es inferior al umbral
configurado ($c < 0.70$), el sistema activa automáticamente
el modo de escalamiento, indicando que el caso debe ser
revisado por un analista humano. Si la confianza supera el
umbral de promoción automática ($c \geq 0.90$), el patrón
detectado se registra en la colección
\texttt{attack\_patterns} de ChromaDB para enriquecer el
\textit{few-shot} de análisis futuros.

La memoria de largo plazo se organiza en tres colecciones
ChromaDB: \texttt{attack\_patterns} (patrones de ataque
confirmados para \textit{few-shot}), \texttt{error\_memory}
(correcciones del analista para evitar errores recurrentes) y
\texttt{host\_history} (historial de diagnósticos validados
por host). Este mecanismo de \textit{feedback loop} permite
que el sistema mejore progresivamente con cada análisis
validado por el analista, aproximándose al comportamiento
acumulativo de un SOC con memoria
institucional~\cite{ismail2025ivam}. En el prototipo
experimental actual, las tres colecciones ChromaDB
contienen un volumen reducido de entradas, adecuado para
la validación del flujo agéntico. La migración hacia bases
de conocimiento más robustas está proyectada para el Proyecto de Título,
donde se evaluará también la integración con el repositorio
oficial MITRE ATT\&CK CTI en formato STIX. El
Algoritmo~\ref{alg:sintesis} describe el proceso completo de
síntesis, generación del reporte y actualización de la
memoria.

\begin{algorithm}[H]
\caption{Síntesis, reporte y feedback loop}
\label{alg:sintesis}
\KwIn{$S$, $R$, diagnóstico $D$, comandos $C_v$,
       confianza $c$}
\KwOut{Reporte $\mathcal{R}$, ChromaDB actualizado}
Construir reporte $\mathcal{R}$ con campos:
    clasificación, confianza, IPs, comandos, validación,
    hipótesis alternativa y escalamiento\;
\eIf{$c < 0.70$}{
    $\mathcal{R}.\texttt{escalamiento} \leftarrow$
        ``Escalar al analista''\;
}{
    $\mathcal{R}.\texttt{escalamiento} \leftarrow$
        ``Auto-resuelto''\;
}
\tcp{Feedback loop: actualizar memoria}
\If{analista valida $\mathcal{R}$}{
    Registrar en \texttt{host\_history}\;
    \If{$c \geq 0.90$}{
        Promover a \texttt{attack\_patterns}\;
    }
}
\If{analista corrige $\mathcal{R}$}{
    Registrar corrección en \texttt{error\_memory}\;
}
\Return{$\mathcal{R}$}\;
\end{algorithm}


%%-------------------------------------------------------------
\section{Componentes de la solución}
%%-------------------------------------------------------------

El agente propuesto está compuesto por nueve componentes funcionales que operan de forma co- ordinada: dos módulos de memoria, cuatro agentes LLM especializados, y tres herramientas de- terminísticas. La interacción entre estos componentes se articula mediante el grafo de estados LangGraph descrito en la sección anterior. La Tabla 8.1 resume los componentes del sistema, su tipo y su función dentro del flujo.

\begin{table}[H]
\centering
\caption{Componentes del sistema agente propuesto.}
\label{tab:componentes}
\renewcommand{\arraystretch}{1.3}
\begin{tabularx}{\textwidth}{p{3.2cm} >{\centering\arraybackslash}p{2.0cm} X}
\toprule
\textbf{Componente} & \textbf{Tipo} & \textbf{Función} \\
\midrule
Memoria de corto plazo  & Memoria  & Traza efímera de un análisis: resumen, plan y observaciones de la sesión actual \\
Memoria de largo plazo  & Memoria  & Persistencia en ChromaDB: patrones de ataque, correcciones y historial por host \\
Agente Triage           & LLM      & Analiza el resumen del lote y decide si requiere contexto histórico (RAG) \\
Agente Selector         & LLM      & Evalúa el diagnóstico y decide si consultar vulnerabilidades CVE en NVD \\
Agente Mitigación       & LLM      & Genera comandos de mitigación Linux; corrige en autorreflexión ante rechazo \\
Agente Auditor Final    & LLM      & Evalúa la calidad global del análisis y asigna nivel de confianza \\
Herramienta RAG         & Tool     & Recupera patrones similares de ChromaDB por similitud coseno \\
Herramienta CVE         & Tool     & Consulta vulnerabilidades conocidas en la API de NVD \\
Herramienta Validador   & Tool     & Verifica comandos de mitigación contra lista blanca de patrones seguros \\
\bottomrule
\end{tabularx}
\end{table}

%%-------------------------------------------------------------
\subsection{Módulos de memoria}
%%-------------------------------------------------------------

El agente propuesto implementa dos niveles de memoria con
responsabilidades distintas. La memoria de corto plazo \texttt{ShortTermMemory} mantiene la traza efímera de un
único análisis: almacena el resumen del lote, el plan generado
por el Planificador y las observaciones acumuladas durante la
ejecución. Esta memoria se reinicia al comenzar cada nuevo
análisis, garantizando que el razonamiento de un lote no
contamine el de un lote posterior.

La memoria de largo plazo (\texttt{LongTermMemory})
persiste entre sesiones mediante ChromaDB con el modelo de
\textit{embeddings} \texttt{all-MiniLM-L6-v2} de Sentence Transformers.
Está organizada en tres colecciones con propósitos
diferenciados. La colección \texttt{attack\_patterns} almacena
patrones de ataque confirmados junto con sus soluciones
validadas, que se recuperan por similitud coseno como
\textit{few-shot} dinámico en la fase de planificación. La
colección \texttt{error\_memory} registra las correcciones
realizadas por el analista ante diagnósticos o comandos
incorrectos, permitiendo que el sistema aprenda de sus propios
errores en análisis futuros. La colección
\texttt{host\_history} mantiene el historial de diagnósticos
validados por host, habilitando la detección de comportamientos
recurrentes o persistentes en un mismo activo.

%%-------------------------------------------------------------
\subsection{Agentes LLM}
%%-------------------------------------------------------------

Los cuatro agentes LLM del sistema utilizan el modelo
\texttt{gemini-2.5-pro} de la API de Google Gemini,
operando con temperatura 0 para maximizar el determinismo de
las respuestas. Cada agente recibe un prompt de sistema
(\textit{system instruction}) que define su rol y el formato
estructurado de salida esperado, lo que facilita el parseado automático entre nodos del grafo.

El \textbf{Agente Triage} actúa como punto de entrada del
grafo. Su instrucción de sistema lo posiciona como analista SOC
experto encargado de interpretar el resumen del lote de logs.
Su primera decisión es si dispone de suficiente contexto
histórico para emitir un diagnóstico preliminar; si no lo
tiene, activa el router hacia la herramienta RAG. Una vez
inyectado el contexto, genera un diagnóstico estructurado con
la clasificación del incidente.

El \textbf{Agente Selector de Herramientas} recibe el
diagnóstico del Triage y decide, entre las herramientas
disponibles (\texttt{query\_cve}, \texttt{validate\_solution}),
cuáles son necesarias para el análisis. Su criterio principal
es si el diagnóstico menciona un servicio específico para el
que puedan existir vulnerabilidades conocidas en NVD. Esta
decisión condicional determina si el flujo pasa por la
herramienta CVE o avanza directamente al agente de mitigación.

El \textbf{Agente Ingeniero de Mitigación} genera comandos
de contención y remediación en Linux (\textit{e.g.},iptables, fail2ban, ufw,
systemctl), tomando en cuenta el diagnóstico, los CVEs
encontrados y las soluciones históricas similares recuperadas
desde ChromaDB. Cuando el Validador rechaza sus comandos, el
agente recibe los motivos de rechazo e implementa un ciclo de
autorreflexión donde corrige los errores identificados, con un
máximo de 2~intentos por análisis.

El \textbf{Agente Auditor Final} evalúa la calidad global del
resultado: la precisión del diagnóstico, la relevancia de los
comandos generados y la evidencia disponible. Su salida es un
nivel de confianza $c \in [0,1]$ que determina el
comportamiento del sistema: si $c < 0.70$, el reporte se
escala al analista humano; si $c \geq 0.90$, el patrón
detectado se promueve automáticamente a la colección
\texttt{attack\_patterns} de ChromaDB.

%%-------------------------------------------------------------
\subsection{Herramientas determinísticas}
%%-------------------------------------------------------------

Las tres herramientas del sistema operan de forma
determinística, sin invocar al modelo de lenguaje, lo que
garantiza su reproducibilidad y trazabilidad.

La \textbf{Herramienta RAG} (\texttt{tool\_consultar\_rag})
consulta la colección \texttt{attack\_patterns} de ChromaDB
mediante la firma del lote actual como vector de consulta.
Recupera hasta $k=3$ patrones similares con un umbral mínimo
de similitud coseno de 0.55, e inyecta los resultados como
contexto al Agente Triage. En el prototipo experimental, la
colección contiene un número reducido de entradas calibradas
para validar el flujo de recuperación; su escala se ampliará
en el Proyecto de Título.

La \textbf{Herramienta CVE} (\texttt{tool\_consultar\_cve})
consulta la API pública de NVD (National Vulnerability
Database) utilizando el servicio identificado en el diagnóstico
como palabra clave de búsqueda. Retorna los CVEs más relevantes
con su descripción y puntuación de severidad, que el Agente de
Mitigación incorpora al contexto de generación de comandos.

La \textbf{Herramienta Validador} (\texttt{tool\_validar\_comandos})
implementa una lista blanca de patrones de comandos Linux
seguros mediante expresiones regulares. Cada comando generado
por el Agente de Mitigación se evalúa contra esta lista; los
comandos que no coincidan son rechazados con un motivo
explícito que se retorna al agente para su corrección. Esta
herramienta actúa como salvaguarda de seguridad del sistema,
impidiendo la propagación de comandos potencialmente
disruptivos o destructivos.

%%-------------------------------------------------------------
\subsection{Routers y aristas condicionales}
%%-------------------------------------------------------------

El grafo LangGraph incorpora tres aristas condicionales
denominadas \textit{routers}, que determinan dinámicamente
el nodo destino según el estado actual del sistema. El
\textbf{Router Triage} evalúa el campo \texttt{necesita\_rag}
del estado: si el Agente Triage no dispone de contexto
histórico, redirige hacia la Herramienta RAG; en caso
contrario, avanza hacia el Agente Selector. El \textbf{Router
Selector} evalúa el campo \texttt{usar\_cve}: si el Agente
Selector determinó que es necesario consultar NVD, redirige
hacia la Herramienta CVE; de lo contrario, avanza directamente
al Agente de Mitigación. El \textbf{Router Validación} evalúa
el campo \texttt{comandos\_validos} y el contador
\texttt{intentos\_correccion}: si los comandos son inválidos y
no se ha alcanzado el máximo de 2~intentos, redirige al Agente
de Mitigación para corrección; en cualquier otro caso, avanza
al Agente Auditor Final.





%============== SIGUEINTE SECCION ============

\section {Tecnologías Consideradas}

La implementación de la arquitectura propuesta requiere integrar diversas tecnologías orientadas al procesamiento de lenguaje natural, la orquestación de agentes inteligentes, la recuperación de conocimiento y el almacenamiento de información. La selección de estas tecnologías se fundamenta tanto en los resultados obtenidos durante la revisión sistemática de la literatura como en los requerimientos funcionales definidos para la solución.

La
Tabla~\ref{tab:stack_tecnologico} presenta el stack
tecnológico completo del sistema.
 
\begin{table}[H]
\centering
\caption{Stack tecnológico del agente propuesto.}
\label{tab:stack_tecnologico}
\renewcommand{\arraystretch}{1.3}
\begin{tabularx}{\textwidth}{p{3.2cm} p{2.2cm} X}
\toprule
\textbf{Tecnología} & \textbf{Versión} &
\textbf{Rol en el sistema} \\
\midrule
Gemini API          & gemini-2.5-pro & Motor LLM de los 4 agentes; temperatura 0, salida JSON \\
google-genai        & SDK oficial           & Cliente Python para la API de Gemini \\
LangGraph           & $\geq$0.2             & Orquestación del grafo de estados (7 nodos, 3 routers) \\
ChromaDB            & $\geq$0.5             & Base de datos vectorial (3 colecciones, modo embebido) \\
Sentence Transformers & all-MiniLM-L6-v2   & Modelo de embeddings (384 dims, similitud coseno) \\
Python              & 3.10+                 & Lenguaje base; parsing, agregación y lógica determinística \\
Requests            & $\geq$2.31            & Cliente HTTP para AbuseIPDB y NVD \\
AbuseIPDB API       & v2                    & Reputación de IPs con caché persistente \\
NVD API             & v2.0 (NIST)           & Consulta de vulnerabilidades CVE por servicio \\
Google Colab        & --                    & Entorno de ejecución del prototipo experimental \\
\bottomrule
\end{tabularx}
\end{table}

Las herramientas consideradas buscan proporcionar capacidades de razonamiento sobre eventos de seguridad, recuperación de contexto desde bases de conocimiento, persistencia de información histórica y generación de reportes estructurados. 

\subsection{Modelos de lenguaje}

Los modelos de lenguaje cumplen un rol central en la arquitectura propuesta, ya que permiten interpretar eventos previamente estructurados, formular hipótesis sobre posibles amenazas, relacionar evidencias y generar una narrativa coherente del incidente. En el prototipo inicial se contempla el uso de un modelo accesible mediante API, con el propósito de validar el funcionamiento general del flujo agentivo y facilitar una implementación experimental de la solución.

Sin embargo, la propuesta no se restringe a un proveedor ni a un modelo específico. Dado que una de las brechas identificadas en la literatura corresponde a la dependencia de modelos propietarios o de gran escala, la arquitectura se concibe como agnóstica al modelo. Esta decisión permite proyectar su evaluación futura con modelos locales de menor tamaño, especialmente en escenarios con restricciones de privacidad, costos, soberanía de datos o disponibilidad de infraestructura.

La Tabla~\ref{tab:modelos_lenguaje} sintetiza las tecnologías consideradas en esta categoría, indicando su rol dentro de la propuesta y la justificación técnica asociada. En particular, se incluyen modelos accesibles mediante API y alternativas locales de menor escala, junto con estrategias de prompting estructurado y generación de salidas en formato JSON, con el fin de facilitar la interoperabilidad entre módulos.

En esta propuesta, los modelos de lenguaje no se emplean como clasificadores aislados, sino como componentes de razonamiento dentro de un flujo controlado. Su función principal es apoyar la interpretación semántica de los eventos, evaluar la evidencia disponible y contribuir a la construcción de una explicación estructurada y trazable del incidente.

\begin{table}[H]
\centering
\caption{Modelos de lenguaje considerados para la propuesta.}
\label{tab:modelos_lenguaje}
\renewcommand{\arraystretch}{1.25}
\begin{tabularx}{\textwidth}{p{2.4cm} X X}
\toprule
\textbf{Tecnología considerada} & \textbf{Rol dentro de la propuesta} & \textbf{Justificación técnica} \\
\midrule
Gemini API & Modelo LLM utilizado para validar el flujo inicial del agente. & Facilita la generación de respuestas estructuradas, el razonamiento sobre eventos y el prototipado rápido de la arquitectura. \\
\midrule
Modelos locales 7B--8B & Alternativa para despliegues con mayor control de datos. & Permiten reducir la dependencia de servicios externos y mejorar la privacidad operacional en entornos con restricciones de seguridad. \\
\midrule
Prompting estructurado & Control del formato de salida del modelo. & Favorece respuestas consistentes, comparables y reutilizables por otros módulos del sistema. \\
\midrule
Salidas en JSON & Estandarización de respuestas del LLM. & Permite integrar el razonamiento del modelo con procesos automáticos de validación, reporte y almacenamiento. \\
\bottomrule
\end{tabularx}
\end{table}

\subsection{Frameworks y bibliotecas principales}

La arquitectura requiere coordinar múltiples etapas de análisis, herramientas externas, mecanismos de memoria y decisiones condicionales. Por ello, se consideran frameworks y bibliotecas que permitan modelar el comportamiento del sistema como un flujo de agentes especializados, en lugar de una ejecución lineal y rígida.

La Tabla~\ref{tab:frameworks_bibliotecas} presenta los principales frameworks y bibliotecas considerados para la implementación de la solución. Entre ellos, LangGraph cumple un rol especialmente relevante, ya que permite representar el proceso de análisis como una máquina de estados compuesta por nodos, transiciones y rutas condicionales. Esta característica se vincula directamente con el objetivo de construir un sistema capaz de ejecutar ciclos de revisión, consulta, validación y refinamiento del diagnóstico.

LangChain se considera como una biblioteca de apoyo para integrar modelos de lenguaje, herramientas externas y mecanismos de memoria. Python, por su parte, actúa como lenguaje base de implementación debido a su flexibilidad para trabajar con APIs, procesamiento de texto, estructuras de datos, validación de comandos y almacenamiento. Finalmente, el uso de JSON permite estandarizar la comunicación entre componentes, lo que favorece la trazabilidad de las entradas y salidas generadas durante el proceso.

\begin{table}[H]
\centering
\caption{Frameworks y bibliotecas considerados para la arquitectura propuesta.}
\label{tab:frameworks_bibliotecas}
\renewcommand{\arraystretch}{1.25}
\begin{tabularx}{\textwidth}{p{2.4cm} X X}
\toprule
\textbf{Framework o biblioteca} & \textbf{Uso considerado} & \textbf{Aporte a la solución} \\
\midrule
LangGraph & Orquestación del flujo agentivo. & Permite representar el análisis como una máquina de estados con nodos, transiciones y decisiones condicionales. \\
\midrule
LangChain & Integración entre LLMs, herramientas y memoria. & Facilita la conexión del modelo con funciones externas, recuperación de conocimiento y procesamiento modular. \\
\midrule
Python & Lenguaje base de implementación. & Entrega flexibilidad para integrar APIs, procesamiento de texto, validación de comandos y almacenamiento. \\
\midrule
Requests / APIs HTTP & Comunicación con servicios externos. & Permite consultar fuentes externas, como NVD o servicios de reputación de direcciones IP. \\
\midrule
JSON & Formato de intercambio entre módulos. & Estandariza entradas y salidas, facilitando la trazabilidad y la validación automática. \\
\bottomrule
\end{tabularx}
\end{table}


\subsection{Herramientas de procesamiento y almacenamiento}

El procesamiento previo de los logs constituye una etapa fundamental de la propuesta, ya que permite transformar registros crudos en evidencia estructurada antes de enviarlos al modelo de lenguaje. Esta fase reduce el ruido, mejora la eficiencia del análisis y disminuye la posibilidad de interpretaciones erróneas. Además, las herramientas consideradas buscan enriquecer el análisis mediante información histórica, inteligencia de amenazas y fuentes externas verificables.

\subsubsection{Preprocesamiento y estructuración de logs}

El preprocesamiento corresponde a la etapa encargada de convertir registros crudos en información estructurada y utilizable por los módulos posteriores de la arquitectura. Esta fase resulta esencial, debido a que los logs de seguridad suelen presentar formatos heterogéneos, información redundante, ruido operativo y eventos distribuidos en distintas fuentes o servicios.

La Tabla~\ref{tab:herramientas_procesamiento} resume las herramientas y técnicas consideradas para esta etapa, entre las que se incluyen expresiones regulares, \textit{parsing }determinístico, agregación de eventos, análisis temporal y validación de comandos. Estas técnicas permiten extraer elementos relevantes, como direcciones IP, usuarios, servicios afectados, códigos de respuesta, eventos críticos, frecuencia de ocurrencia y ventanas temporales.

El uso de \textit{parsing }determinístico permite normalizar los logs antes de enviarlos al modelo de lenguaje. Esta decisión reduce la complejidad de la entrada y contribuye a disminuir el riesgo de interpretaciones ambiguas. A su vez, la agregación de eventos facilita la identificación de patrones repetitivos, como intentos de fuerza bruta, escaneo de servicios, accesos fallidos o comportamientos anómalos asociados a una misma entidad.

El análisis temporal cumple un rol importante en la reconstrucción del incidente, ya que permite ordenar los eventos según su ocurrencia y establecer posibles relaciones entre ellos. Esta capacidad es especialmente relevante para diferenciar eventos aislados de secuencias de ataque más amplias. Por último, la validación de comandos se incorpora como una medida de control sobre las acciones de mitigación generadas, con el propósito de reducir recomendaciones inseguras, ambiguas o potencialmente disruptivas.

En conjunto, estas herramientas proporcionan una base estructurada y verificable para el razonamiento posterior del agente. De este modo, el modelo de lenguaje no interpreta directamente registros sin procesar, sino una representación más clara, contextualizada y trazable de la evidencia disponible.

\begin{table}[H]
\centering
\caption{Herramientas y técnicas de procesamiento consideradas.}
\label{tab:herramientas_procesamiento}
\renewcommand{\arraystretch}{1.25}
\begin{tabularx}{\textwidth}{p{2.4cm} X X}
\toprule
\textbf{Herramienta o técnica} & \textbf{Función dentro del sistema} & \textbf{Resultado esperado} \\
\midrule
Expresiones regulares & Extracción de patrones desde logs. & Identificación de direcciones IP, usuarios, servicios, eventos críticos y códigos de respuesta. \\
\midrule
Parsing determinístico & Normalización inicial de eventos. & Conversión de logs crudos en estructuras analizables por los módulos posteriores. \\
\midrule
Agregación de eventos & Agrupación por IP, usuario, servicio o ventana temporal. & Detección de patrones repetitivos, como intentos de fuerza bruta o escaneo. \\
\midrule
Análisis temporal & Ordenamiento de eventos según su ocurrencia. & Reconstrucción de la secuencia del incidente. \\
\midrule
Validación de comandos & Revisión de acciones de mitigación generadas. & Reducción de recomendaciones inseguras, ambiguas o potencialmente disruptivas. \\
\bottomrule
\end{tabularx}
\end{table}

\subsubsection{Almacenamiento y recuperación de conocimiento}

Esta sección corresponde a las herramientas destinadas a almacenar, recuperar y enriquecer el conocimiento utilizado durante el análisis. Su propósito es complementar la inferencia del modelo de lenguaje con información previamente registrada o proveniente de fuentes externas especializadas.

La Tabla~\ref{tab:almacenamiento_recuperacion} presenta las principales tecnologías y fuentes consideradas en esta área. ChromaDB se contempla como base vectorial para almacenar casos previos, patrones de ataque, fragmentos de conocimiento y representaciones semánticas reutilizables. Esta funcionalidad permite que el agente compare los eventos actuales con incidentes históricos o con conocimiento incorporado previamente.

MITRE ATT\&CK CTI se considera una fuente estructurada de inteligencia de amenazas, orientada a vincular los eventos observados con tácticas, técnicas, subtécnicas, procedimientos y posibles grupos de amenaza. Su integración mediante STIX permite trabajar con un formato estandarizado, lo que facilita la incorporación de este conocimiento en una base consultable por el sistema.

Además, se consideran fuentes externas como NVD y AbuseIPDB. NVD permite consultar vulnerabilidades conocidas asociadas a servicios, productos o tecnologías identificadas durante el análisis. AbuseIPDB, en cambio, proporciona información de reputación sobre direcciones IP involucradas en los eventos. Estas fuentes aportan evidencia complementaria para fortalecer o debilitar las hipótesis generadas por el agente.

Esta área resulta relevante porque evita que el diagnóstico dependa exclusivamente del conocimiento interno del modelo de lenguaje. Al incorporar recuperación semántica y fuentes externas, la arquitectura puede fundamentar mejor sus conclusiones, mejorar la trazabilidad del razonamiento y entregar reportes más útiles para el analista.

\begin{table}[H]
\centering
\caption{Herramientas de almacenamiento y recuperación de conocimiento consideradas.}
\label{tab:almacenamiento_recuperacion}
\renewcommand{\arraystretch}{1.25}
\begin{tabularx}{\textwidth}{p{2.4cm} X X}
\toprule
\textbf{Tecnología o fuente} & \textbf{Uso considerado} & \textbf{Aporte a la propuesta} \\
\midrule
ChromaDB & Base vectorial para memoria y recuperación semántica. & Permite almacenar casos previos, patrones de ataque y conocimiento estructurado. \\
\midrule
MITRE ATT\&CK CTI & Base de conocimiento sobre tácticas, técnicas y procedimientos. & Permite asociar eventos observados con técnicas reconocidas de ataque. \\
\midrule
STIX & Formato estructurado para inteligencia de amenazas. & Facilita poblar la base de conocimiento desde fuentes estandarizadas. \\
\midrule
NVD & Consulta de vulnerabilidades CVE. & Permite relacionar servicios o tecnologías afectadas con vulnerabilidades conocidas. \\
\midrule
AbuseIPDB & Consulta de reputación de direcciones IP. & Aporta evidencia externa sobre direcciones IP potencialmente maliciosas. \\
\bottomrule
\end{tabularx}
\end{table}

\subsection{Síntesis tecnológica de la propuesta}

La selección de tecnologías descrita permite estructurar la solución en capas funcionales, cada una con un propósito específico dentro del flujo de análisis: recibir y preparar los logs, razonar sobre los eventos, coordinar agentes y herramientas, y enriquecer el diagnóstico mediante conocimiento externo.

La Tabla~\ref{tab:sintesis_tecnologica} resume esta organización por capas. La capa de entrada y procesamiento se orienta a transformar logs crudos en evidencia estructurada. La capa de razonamiento emplea modelos de lenguaje para interpretar eventos y formular hipótesis. La capa de orquestación agentiva coordina los módulos del sistema mediante un flujo basado en decisiones condicionales. Por último, la capa de conocimiento y validación incorpora memoria, inteligencia de amenazas, consultas externas y mecanismos de control sobre las respuestas generadas.

En conjunto, estas tecnologías permiten proponer una arquitectura que no depende exclusivamente de un modelo de lenguaje, sino de la integración entre procesamiento determinístico, razonamiento semántico, recuperación de conocimiento y validación de acciones. Esta combinación resulta coherente con el objetivo de avanzar desde el resumen automático de logs hacia una reconstrucción más completa, trazable y accionable de incidentes de seguridad.

\begin{table}[H]
\centering
\caption{Síntesis de las capas tecnológicas consideradas en la propuesta.}
\label{tab:sintesis_tecnologica}
\renewcommand{\arraystretch}{1.25}
\begin{tabularx}{\textwidth}{p{2.4cm} X X}
\toprule
\textbf{Capa} & \textbf{Tecnologías asociadas} & \textbf{Propósito} \\
\midrule
Capa de entrada y procesamiento & Python, expresiones regulares, parsing y agregación temporal. & Estructurar logs crudos y extraer evidencia relevante para el análisis. \\
\midrule
Capa de razonamiento & Gemini API, modelos locales, prompting estructurado y JSON. & Interpretar eventos, formular hipótesis y generar explicaciones estructuradas. \\
\midrule
Capa de orquestación agentiva & LangGraph y LangChain. & Coordinar agentes, herramientas, memoria y decisiones condicionales dentro del flujo de análisis. \\
\midrule
Capa de conocimiento y validación & ChromaDB, MITRE ATT\&CK, STIX, NVD, AbuseIPDB y validación de comandos. & Enriquecer el análisis, recuperar evidencia externa y controlar la seguridad de las respuestas generadas. \\
\bottomrule
\end{tabularx}
\end{table}



%%-------------------------------------------------------------
\section{Requerimientos técnicos}
%%-------------------------------------------------------------

Esta sección especifica los requerimientos de hardware,
software y servicios externos necesarios para ejecutar el
prototipo experimental del agente propuesto. Los
requerimientos se organizan en tres categorías: entorno de
ejecución, dependencias de software y acceso a servicios
externos.

%%-------------------------------------------------------------
\subsection{Requerimientos de hardware}
%%-------------------------------------------------------------

El prototipo experimental se ejecuta sobre Google Colab, que
provisiona los recursos de cómputo de forma dinámica sin
requerir hardware dedicado. En un entorno de ejecución local
alternativo, los requerimientos mínimos son los siguientes.

\begin{table}[H]
\centering
\caption{Requerimientos de hardware para ejecución local
         del prototipo.}
\label{tab:req_hardware}
\renewcommand{\arraystretch}{1.3}
\begin{tabularx}{\textwidth}{p{3.0cm} p{3.5cm} X}
\toprule
\textbf{Componente} & \textbf{Mínimo} & \textbf{Recomendado} \\
\midrule
Procesador  & CPU de 4 núcleos & CPU de 8 núcleos \\
Memoria RAM & 8 GB             & 16 GB \\
Almacenamiento & 2 GB libres   & 10 GB libres (para ChromaDB y caché de modelos) \\
Conectividad & Internet estable & Conexión de baja latencia para APIs Gemini, AbuseIPDB y NVD \\
\bottomrule
\end{tabularx}
\end{table}


%%-------------------------------------------------------------
\subsection{Requerimientos de software}
%%-------------------------------------------------------------

El sistema requiere Python 3.10 o superior y las siguientes
dependencias, instalables mediante \texttt{pip}:

\begin{table}[H]
\centering
\caption{Dependencias de software del prototipo.}
\label{tab:req_software}
\renewcommand{\arraystretch}{1.3}
\begin{tabularx}{\textwidth}{>{\raggedright\arraybackslash}p{4.8cm} >{\centering\arraybackslash}p{1.6cm} X}
\toprule
\textbf{Paquete} & \textbf{Versión mín.} &
\textbf{Propósito} \\
\midrule
\texttt{google-genai}         &  1.0  & Cliente oficial SDK para la API de Google Gemini \\
\texttt{chromadb}             &  0.5  & Base de datos vectorial para memoria de largo plazo \\
\texttt{sentence-transformers}&  2.7  & Modelo de embeddings \texttt{all-MiniLM-L6-v2} \\
\texttt{langgraph}            &  0.2  & Orquestación del grafo de estados agéntico \\
\texttt{requests}             &  2.31 & Cliente HTTP para APIs externas (AbuseIPDB, NVD) \\
\texttt{typing-extensions}    &  4.0  & Soporte para \texttt{TypedDict} en Python 3.10 \\
\bottomrule
\end{tabularx}
\end{table}


%%-------------------------------------------------------------
\subsection{Requerimientos de servicios externos}
%%-------------------------------------------------------------

El prototipo depende de tres servicios externos accesibles
mediante HTTPS, cada uno con sus propios requerimientos de
credenciales y cuota de uso.

\begin{table}[H]
\centering
\caption{Requerimientos de acceso a servicios externos.}
\label{tab:req_servicios}
\renewcommand{\arraystretch}{1.3}
\begin{tabularx}{\textwidth}{p{2.8cm} p{2.2cm} p{2.2cm} X}
\toprule
\textbf{Servicio} & \textbf{Autenticación} &
\textbf{Cuota (plan gratuito)} & \textbf{Uso en el sistema} \\
\midrule
Google Gemini API &
    API Key requerida &
    60 req/min &
    Motor LLM de los 4 agentes \\
AbuseIPDB API &
    API Key requerida &
    1.000 consultas/día &
    Reputación de IPs (con caché para minimizar uso) \\
NVD API (NIST) &
    Sin autenticación (pública) &
    5 req/30 s sin key; 50 req/30 s con key &
    Consulta de CVEs por servicio afectado \\
\bottomrule
\end{tabularx}
\end{table}

Las credenciales de la API de Gemini y AbuseIPDB se
gestionan mediante el sistema de secretos de Google Colab
(\texttt{userdata.get()}) o variables de entorno del sistema
operativo, sin almacenarlas en el código fuente. El módulo
\texttt{ReputationCache} garantiza que la cuota de AbuseIPDB
se utilice eficientemente: cada IP distinta se consulta
exactamente una vez por sesión, sirviendo las consultas
subsecuentes desde el archivo \texttt{ip\_reputation\_cache.json}.

%%-------------------------------------------------------------
\subsection{Datasets de evaluación}
%%-------------------------------------------------------------

Para la validación experimental del agente propuesto se
consideran cuatro datasets de referencia ampliamente utilizados
en la literatura de análisis de logs y detección de anomalías.
La selección se fundamenta en su presencia en el corpus SLR y en su relevancia para el dominio de logs de
seguridad que el agente procesa.

\begin{table}[H]
\centering
\caption{Datasets considerados para la evaluación del
         agente propuesto.}
\label{tab:datasets}
\renewcommand{\arraystretch}{1.3}
\begin{tabularx}{\textwidth}{p{3.2 cm} p{2.8cm} p{1.8cm} X}
\toprule
\textbf{Dataset} & \textbf{Tipo de logs} &
\textbf{Tamaño aprox.} & \textbf{Relevancia para el agente} \\
\midrule
HDFS%~\cite{pham2024multilog} 
&
    Sistema distribuido (Hadoop) &
    11M líneas &
    Benchmark estándar con ground truth etiquetado; permite
    evaluar la detección de anomalías en logs de sistema \\
BGL%~\cite{ali2025mlstudy} 
&
    Supercomputador (kernel) &
    5M líneas &
    Logs de sistema similares a los que procesa el agente;
    incluye eventos de fallo y comportamiento anómalo \\
UWF-ZeekDataFall22%~\cite{krebs2024svm} 
&
    Red (Zeek/IDS) con tácticas MITRE &
    700K registros &
    Dataset de red con etiquetas MITRE ATT\&CK; directamente
    alineado con el dominio de análisis de incidentes del
    agente \\
LogHub%~\cite{ma2024llmparser} 
&
    Colección de 16 sistemas &
    Varios &
    Colección de referencia para evaluación de parsers;
    permite comparar la etapa de preprocesamiento del agente
    con los parsers del corpus \\
\bottomrule
\end{tabularx}
\end{table}

En el prototipo experimental, el
agente se valida sobre muestras representativas de estos
datasets, con el objetivo de verificar el correcto funcionamiento
del flujo agéntico completo \textit{i.e.},  preprocesamiento, planificación,
ejecución y síntesis; sobre logs de distintos formatos y
dominios. La evaluación exhaustiva sobre los datasets
completos, con métricas de precisión, coherencia narrativa
y comparación con baselines, está proyectada para el
Proyecto de Título.

%===============================
%%-------------------------------------------------------------
\section{Diseño Experimental y Metodología de Evaluación}
%%-------------------------------------------------------------

Para validar la viabilidad técnica y la efectividad de la arquitectura del agente propuesto, se llevará a cabo una evaluación experimental basada en métodos cuantitativos y cualitativos. El objetivo es determinar la capacidad del agente para procesar logs heterogéneos, identificar incidentes de seguridad con precisión y generar narrativas coherentes y accionables, comparando su desempeño contra enfoques tradicionales.

Esta sección detalla el diseño de los experimentos controlados, las métricas de desempeño seleccionadas y el procedimiento de comparación con \textit{baselines} basados en reglas.

\subsection{Diseño de experimentos controlados}

La evaluación se estructurará en torno a experimentos controlados utilizando muestras representativas de los datasets descritos anteriormente en la Tabla~\ref{tab:datasets}. Dado el carácter prototípico de esta etapa del seminario, no se procesarán los datasets en su totalidad, sino lotes \textit{i.e., batches}, de logs seleccionados que contengan tanto comportamiento normal como anomalías de seguridad inyectadas o documentadas en el \textit{ground truth} de cada dataset.

Se definen dos escenarios experimentales principales:

\begin{enumerate}
    \item \textbf{Escenario A: Detección e interpretación de anomalías de sistema:} Utilizando muestras de HDFS y BGL, se evaluará la capacidad del agente para identificar eventos de falla o accesos inómalos que, aunque no siempre constituyen un ataque dirigido, representan riesgos operativos que requieren diagnóstico semántico.
    \item \textbf{Escenario B: Análisis de incidentes de red con alineación MITRE:} Utilizando el dataset UWF-ZeekDataFall22, se probará la capacidad del flujo agéntico completo para correlacionar alertas de red, consultar fuentes externas (NVD, AbuseIPDB), asociar tácticas/técnicas MITRE ATT\&CK y generar recomendaciones de mitigación válidas.
\end{enumerate}

\subsection{Métricas de desempeño}

Para evaluar el cumplimiento de los objetivos, se han seleccionado métricas que abordan tanto la exactitud técnica de la detección como la calidad lingüística y operativa de la salida generada por el LLM.

\subsubsection{Métricas cuantitativas de identificación de incidentes}

Se utilizarán métricas estándar de clasificación binaria para evaluar la capacidad del sistema para diferenciar entre actividad normal y actividad maliciosa/anómala, basándose en la clasificación final del reporte generado ($c \geq 0.70$ para incidentes confirmados, $c < 0.70$ o falsos positivos para escalamiento/error).

\begin{itemize}
    \item \textbf{Precisión (\textit{Precision}):} Proporción de incidentes reportados por el agente que realmente corresponden a anomalías según el \textit{ground truth}. Es crítica para minimizar la fatiga de alertas del analista.
    \item \textbf{Exhaustividad (\textit{Recall}):} Proporción de anomalías reales presentes en el dataset que fueron correctamente identificadas por el agente. Es crítica para no pasar por alto amenazas reales.
    \item \textbf{F1-Score:} Media armónica entre Precisión y Exhaustividad, proporcionando una métrica balanceada del desempeño de detección.
\end{itemize}

\subsection{Métricas cualitativas de la narrativa operativa}

Dado que una salida estructurada en JSON puede ser sintácticamente correcta pero semánticamente inútil, se empleará una metodología de evaluación cualitativa basada en rúbricas de expertos (humanos) para evaluar los campos de texto libre del reporte final (\texttt{justificacion}, \texttt{hipotesis\_alternativa}). Se evaluarán tres dimensiones:

\begin{itemize}
    \item \textbf{Coherencia narrativa:} Se evaluará si el diagnóstico generado sigue una lógica consistente con la evidencia presentada en el resumen de logs $S$. Se utilizará una escala de Likert de 1 a 5 (1: Totalmente incoherente, 5: Totalmente coherente).
    \item \textbf{Accionabilidad de la mitigación:} Se evaluará si los comandos validados $C_v$ y las recomendaciones son técnicamente correctos, pertinentes para el incidente detectado y seguros de aplicar. Escala de 1 a 5 (1: Inútil/Peligroso, 5: Directamente aplicable).
    \item \textbf{Trazabilidad del razonamiento:} Se verificará si el reporte explicita claramente cómo se utilizó el contexto recuperado (RAG o CVE) para llegar al diagnóstico final. Escala Binaria (Sí/No).
\end{itemize}

\subsection{Comparación con enfoques tradicionales}

Para establecer un \textit{baseline} de comparación, los mismos lotes de logs experimentales serán procesados mediante un enfoque tradicional basado en reglas determinísticas (\textit{e.g.}, patrones de coincidencia de cadenas, umbrales de frecuencia simples).

La comparación se centrará en demostrar el valor agregado del análisis semántico del LogGuard Agent frente a las reglas rígidas:

\begin{itemize}
    \item \textbf{Capacidad de generalización:} Se espera demostrar que el agente puede identificar variantes de ataques que no coinciden exactamente con una regla prefijada, gracias al razonamiento semántico del LLM.
    \item \textbf{Reducción de falsos positivos:} Se evaluará si el agente, al considerar el contexto completo del lote y fuentes externas, logra descartar alertas ruidosas que un sistema basado en reglas simples clasificaría erróneamente como críticas.
    \item \textbf{Riqueza de la información:} Se comparará la salida unidimensional de una regla (ej. "Alerta: Fuerza BrutaDetectada") frente al reporte estructurado y multidimensional del agente, evaluando cuál aporta mayor valor operativo al analista humano.
\end{itemize}


%===============================
\section{Alcance de la Implementación}

El alcance de la implementación se define en función del carácter prototípico de la solución propuesta. En esta etapa, el objetivo no es desarrollar una plataforma SOC completa ni un sistema listo para operar en ambientes productivos, sino construir una arquitectura experimental que permita validar la viabilidad técnica de un agente LLM para el análisis semántico de logs de seguridad. Por lo tanto, la implementación se orienta a demostrar el flujo principal de análisis, la integración de componentes clave y la generación de una salida estructurada que apoye la reconstrucción de incidentes.

\subsection{Capacidades que se espera implementar}

La implementación contempla, en primer lugar, un módulo de ingreso y preprocesamiento de logs. Este componente tendrá como propósito recibir un conjunto acotado de registros de seguridad, normalizar su contenido y extraer elementos relevantes para el análisis posterior. Entre estos elementos se consideran direcciones IP, usuarios, servicios afectados, eventos críticos, códigos de respuesta, frecuencia de ocurrencia y ventanas temporales. Esta etapa permitirá transformar logs crudos en una representación estructurada, reduciendo el ruido de entrada y facilitando el razonamiento del agente.

En segundo lugar, se espera implementar un flujo agentivo de análisis mediante LangGraph. Este flujo estará compuesto por nodos especializados, encargados de realizar tareas diferenciadas dentro del proceso de análisis. Entre estas tareas se incluyen el análisis inicial del evento, la decisión sobre la necesidad de consultar memoria o conocimiento externo, la recuperación de información relevante, la generación de hipótesis, la propuesta de acciones de mitigación y la evaluación final del resultado. Esta organización permitirá representar el análisis como un proceso iterativo, en lugar de un pipeline lineal de entrada y salida.

Una tercera capacidad corresponde a la incorporación de recuperación aumentada por conocimiento. Para ello, se considera el uso de una base vectorial, como ChromaDB, destinada a almacenar información histórica, patrones de ataque y conocimiento estructurado. En una versión extendida de la propuesta, esta base podrá incorporar información proveniente de MITRE ATT\&CK en formato STIX, con el objetivo de asociar los eventos observados con tácticas, técnicas y procedimientos reconocidos en ciberseguridad.

Asimismo, se contempla la integración con fuentes externas de apoyo al análisis. Entre ellas se consideran servicios de consulta de vulnerabilidades, como NVD, y fuentes de reputación de direcciones IP, como AbuseIPDB. Estas herramientas permitirán complementar la interpretación del modelo de lenguaje con evidencia externa, fortaleciendo la trazabilidad del diagnóstico y reduciendo la dependencia exclusiva del conocimiento interno del modelo.

Otra capacidad esperada es la generación de acciones de mitigación o contención. Estas acciones estarán orientadas a entregar recomendaciones técnicas frente a los hallazgos identificados, tales como bloqueo de direcciones IP, revisión de usuarios comprometidos, endurecimiento de servicios o verificación de configuraciones. Sin embargo, estas acciones no serán ejecutadas automáticamente, sino que serán presentadas como recomendaciones sujetas a revisión.

Adicionalmente, se espera implementar un módulo básico de validación de comandos o acciones sugeridas. Este componente tendrá como finalidad identificar instrucciones potencialmente inseguras, ambiguas o disruptivas antes de incorporarlas al reporte final. Con ello se busca introducir una capa mínima de control sobre las respuestas generadas por el modelo, especialmente cuando estas incluyen comandos o medidas operativas.

Finalmente, la implementación deberá generar un reporte estructurado del incidente. Este reporte incluirá, al menos, una descripción general del evento, la secuencia temporal de los registros relevantes, las entidades involucradas, la hipótesis del incidente, las evidencias utilizadas, las posibles técnicas MITRE ATT\&CK asociadas, el nivel de confianza del análisis y las acciones recomendadas. Esta salida constituye uno de los elementos centrales de la propuesta, ya que permite avanzar desde la simple detección de eventos hacia una reconstrucción comprensible y accionable del incidente.

\subsection{Limitaciones conocidas y aspectos fuera del alcance}

Dado el carácter experimental de la propuesta, la implementación tendrá un alcance acotado. En primer lugar, no se considera el desarrollo de una solución SOC productiva ni su integración completa con plataformas SIEM reales. Aunque la arquitectura propuesta podría proyectarse hacia herramientas de este tipo, la implementación se limitará a validar el flujo de análisis sobre conjuntos de logs controlados o previamente definidos.

En segundo lugar, no se contempla la ingesta masiva de logs en tiempo real ni la evaluación de desempeño en ambientes de alta concurrencia. El prototipo se enfocará en el análisis semántico y en la reconstrucción del incidente, por lo que aspectos como escalabilidad, tolerancia a fallos, balanceo de carga, almacenamiento distribuido o procesamiento continuo quedarán fuera del alcance de esta etapa.

Tampoco se considera la ejecución automática de acciones de respuesta sobre sistemas reales. Las mitigaciones generadas por el agente serán tratadas como recomendaciones técnicas y deberán ser revisadas por un analista humano antes de cualquier aplicación. Esta delimitación es necesaria para evitar riesgos operacionales, especialmente en acciones que puedan afectar disponibilidad de servicios, conectividad de red o cuentas de usuario.

Otra limitación corresponde al tipo y diversidad de logs utilizados. La implementación se centrará en fuentes representativas para el análisis de incidentes, tales como registros de autenticación, eventos de servicios web o mensajes de sistema. No se busca cubrir todos los formatos existentes ni desarrollar un parser universal para cualquier tecnología. La normalización de logs se abordará mediante reglas y patrones definidos para los casos de prueba seleccionados.

Asimismo, la propuesta no contempla el entrenamiento o ajuste fino de modelos de lenguaje. El prototipo utilizará modelos previamente disponibles, mediante API o eventualmente mediante ejecución local, y se concentrará en la orquestación, el prompting estructurado, la recuperación de conocimiento y la validación del flujo. La comparación exhaustiva entre distintos modelos, tamaños o estrategias de fine-tuning podrá considerarse como trabajo futuro.

La integración con MITRE ATT\&CK se plantea como parte del diseño de la arquitectura y, según el avance de implementación, podrá abordarse de forma parcial. En esta etapa, el objetivo principal será demostrar la utilidad de asociar eventos con tácticas y técnicas conocidas, sin pretender cubrir exhaustivamente la totalidad de la base ATT\&CK ni garantizar una atribución completa de grupos de amenaza.

También quedan fuera del alcance tareas especializadas como análisis forense digital profundo, inspección de tráfico de red a nivel de paquetes, análisis de malware, correlación con telemetría EDR, generación automática de reglas SIEM o evaluación formal de impacto sobre infraestructura real. Estas funciones exceden el objetivo del prototipo y requerirían componentes adicionales no contemplados en esta etapa.

Por último, tal como se mencionó previamente, la solución no busca reemplazar al analista de seguridad. Su propósito es actuar como un sistema de apoyo para organizar evidencia, sugerir hipótesis, recuperar conocimiento relevante y producir una narrativa inicial del incidente. La validación final, la priorización de acciones y la toma de decisiones operacionales permanecerán bajo completa responsabilidad humana.

En síntesis, el alcance de la implementación se concentra en validar una arquitectura modular y agentiva para el análisis de logs de seguridad. La propuesta busca demostrar que es posible integrar preprocesamiento, razonamiento con LLMs, recuperación de conocimiento, consulta de fuentes externas, validación de acciones y generación de reportes en un flujo coherente. Sus limitaciones se relacionan principalmente con el carácter experimental del prototipo, la ausencia de despliegue productivo y la necesidad de supervisión humana en las decisiones de respuesta.
